\subsubsection{Oracle Action and Phase Kickback}\label{ex:oracle-action-and-phase-kickback}

This section is the heart of Deutsch's algorithm, even we spoiled it in the previous section. Up to now we've prepared the \emph{right} input state. Here we show \textbf{why that preparation was useful}: the oracle no longer stores the function value in the second qubit; it \textbf{kicks it back as a phase} onto the first qubit.

\highspace
In the previous section we prepared the input state:
\begin{equation*}
    \ket{\psi_1} = \ket{+} \ket{-} = \frac{\ket{0} + \ket{1}}{\sqrt{2}} \otimes \ket{-}
\end{equation*}
Now we apply the oracle:
\begin{equation*}
    U_f \ket{x, y} = \ket{x, y \oplus f(x)}
\end{equation*}
At first sight, it may seem that the oracle should modify only the computation qubit, since this is the qubit on which the XOR operation is performed. Surprisingly, once the computation qubit is prepared in the state $\ket{-}$, the oracle leaves it unchanged and instead transfers the information about $f(x)$ into the phase of the data qubit. This phenomenon is known as \textbf{phase kickback}.

\highspace
\begin{flushleft}
    \textcolor{Green3}{\faIcon{tools} \textbf{Oracle Action on the State $\ket{-}$}}
\end{flushleft}
Recall that:
\begin{equation*}
    \ket{-} = \frac{\ket{0} - \ket{1}}{\sqrt{2}}
\end{equation*}
And that the oracle acts as:
\begin{equation*}
    U_f \ket{x,y} = \ket{x, y \oplus f(x)}
\end{equation*}
Let us evaluate the two possible cases:
\begin{itemize}
    \item \hl{Case $f(x) = 0$.} The oracle performs no bit flip:
    \begin{equation*}
        U_f \ket{x}\ket{-} = \ket{x}\ket{-}
    \end{equation*}
    Nothing changes, and the computation qubit remains in the state $\ket{-}$.

    
    \item \hl{Case $f(x) = 1$.} The oracle flips the computation qubit:
    \begin{equation*}
        U_f \ket{x}\ket{-} = \ket{x} X\ket{-}
    \end{equation*}
    Since the Pauli-$X$ operator has eigenvalue $-1$ on the state $\ket{-}$, we have:
    \begin{equation*}
        X\ket{-} = -\ket{-}
        \quad \implies \quad
        U_f \ket{x}\ket{-} = -\ket{x}\ket{-}
    \end{equation*}
    The computation qubit is still in the state \ket{-}; only a minus sign has appeared.
\end{itemize}
Both cases can therefore be written compactly as:
\begin{equation}
    U_f \ket{x}\ket{-} = (-1)^{f(x)} \ket{x}\ket{-} \quad \text{where} \quad \left(-1\right)^{f(x)} = \begin{cases}
        +1 & \text{if } f(x) = 0 \\
        -1 & \text{if } f(x) = 1
    \end{cases}
\end{equation}
This equation is the key identity behind Deutsch's algorithm.

\highspace
\begin{flushleft}
    \textcolor{Green3}{\faIcon{tools} \textbf{Applying the Oracle to the Superposition}}
    \phantomsection\label{sec:deutsch-applying-the-oracle-to-the-superposition}
\end{flushleft}
The data qubit is not in a basis state but in the superposition:
\begin{equation*}
    \ket{+} = \frac{\ket{0} + \ket{1}}{\sqrt{2}}
\end{equation*}
Therefore, the input state to the oracle is:
\begin{equation*}
    \ket{\psi_1} = \ket{+}\ket{-} = \frac{\ket{0} + \ket{1}}{\sqrt{2}} \otimes \ket{-}
\end{equation*}
The oracle acts on each branch of the superposition independently:
\begin{align*}
    U_f \ket{\psi_1} &= \frac{1}{\sqrt{2}} \left( U_f \ket{0, -} + U_f \ket{1, -} \right) \\
    &= \frac{1}{\sqrt{2}} \left( (-1)^{f(0)} \ket{0, -} + (-1)^{f(1)} \ket{1, -} \right)
\end{align*}
Factoring out the unchanged computation qubit ($\ket{-}$), we have:
\begin{equation}
    U_f \ket{\psi_1} = \frac{1}{\sqrt{2}} \left( (-1)^{f(0)} \ket{0} + (-1)^{f(1)} \ket{1} \right) \otimes \ket{-}
\end{equation}
Notice that:
\begin{itemize}
    \item The \textbf{computation qubit} is exactly the \textbf{same as before} ($\ket{-}$);
    \item All the \textbf{information} about the function has been \textbf{transferred into the relative signs of the data qubit}.
\end{itemize}
This is precisely the \textbf{phase kickback} mechanism!

\highspace
The oracle \textbf{no} longer \textbf{stores the values} of the function as bits (as in the classical case). Instead, it \textbf{encodes} them as \textbf{phase factors}:
\begin{equation*}
    \left(-1\right)^{f(0)} \quad \text{and} \quad \left(-1\right)^{f(1)}
\end{equation*}
Each branch of the superposition acquires its own sign:
\begin{equation*}
    \begin{aligned}
        \ket{0} &\longrightarrow (-1)^{f(0)} \ket{0} \\
        \ket{1} &\longrightarrow (-1)^{f(1)} \ket{1}
    \end{aligned}
\end{equation*}
The quantum state therefore contains both function evaluations simultaneously:
\begin{equation*}
    \frac{1}{\sqrt{2}} \left( (-1)^{f(0)} \ket{0} + (-1)^{f(1)} \ket{1} \right)
\end{equation*}
Not as two readable output bits, but as a \textbf{pattern of relative phases}. These relative phases are exactly what the final Hadamard gate will transform into a measurable outcome.

\newpage

\begin{flushleft}
    \textcolor{Green3}{\faIcon{tags} \textbf{Constant and Balanced Phase Patterns}}
\end{flushleft}
The four possible functions produce four different phase patterns:
\begin{center}
    \begin{tabular}{@{} c c c c @{}}
        \toprule
        \textbf{Function} & $f(0)$ & $f(1)$ & \textbf{Data qubit after the oracle} \\
        \midrule
        $f(x)=0$            & 0 & 0 & $\dfrac{\ket{0} + \ket{1}}{\sqrt{2}}=\ket{+}$ \\[1.2em]
        $f(x)=1$            & 1 & 1 & $-\dfrac{\ket{0} + \ket{1}}{\sqrt{2}}=-\ket{+}$ \\
        \cmidrule{1-4}
        $f(x)=x$            & 0 & 1 & $\dfrac{\ket{0} - \ket{1}}{\sqrt{2}}=\ket{-}$ \\[1.2em]
        $f(x)=\overline{x}$ & 1 & 0 & $-\dfrac{\ket{0} - \ket{1}}{\sqrt{2}}=-\ket{-}$ \\
        \bottomrule
    \end{tabular}
\end{center}

\noindent
A remarkable pattern immediately appears:
\begin{itemize}
    \item \important{Constant functions} produce:
    \begin{equation*}
        \pm \ket{+}
    \end{equation*}
    Which differ only by a \textbf{global phase}.
    \item \important{Balanced functions} produce:
    \begin{equation*}
        \pm \ket{-}
    \end{equation*}
    Which again differ only by a \textbf{global phase}.
\end{itemize}
Since global phases are physically unobservable, the oracle has effectively reduced the four possible functions to \textbf{two distinguishable quantum states}:
\begin{equation*}
    \text{Constant} \longrightarrow \ket{+} \qquad \text{and} \qquad \text{Balanced} \longrightarrow \ket{-}
\end{equation*}
\hl{This is exactly why the algorithm can determine whether the function is constant or balanced with a single oracle query.} The only remaining task is to convert the states \ket{+} and \ket{-} into computational basis states, which is accomplished by the final Hadamard gate (next section).